\documentclass[11pt,a4paper]{article}

\usepackage[T1]{fontenc}
\usepackage[utf8]{inputenc}
\usepackage[italian]{babel}
\usepackage{amsmath,amssymb}
\usepackage{booktabs}
\usepackage{geometry}
\usepackage[protrusion=true,expansion=false]{microtype}
\usepackage[colorlinks=true,linkcolor=black,urlcolor=blue]{hyperref}

\geometry{margin=2.6cm}

% --- abbreviazioni unità ---
\newcommand{\un}[1]{\,\mathrm{#1}}

\title{Calcolo dell'induttanza di dispersione\\di un trasformatore trifase dai dati di targa}
\author{Nota tecnica}
\date{\today}

\begin{document}
\maketitle

\begin{abstract}
\noindent
Il documento ricava l'induttanza di dispersione (leakage inductance) di un trasformatore
trifase a partire dai soli dati di targa $S_n$ (o $P_n$), $U_n$, $I_n$, $u_{cc}$, con
l'aggiunta opzionale delle perdite di cortocircuito $P_{cc}$. Vengono trattati: il sistema
per unità, la separazione fra parte resistiva e reattiva, la ripartizione fra primario e
secondario, il riporto fra i due lati, le insidie legate al collegamento
stella/triangolo, una verifica geometrica indipendente con la formula di Rogowski e una
serie di note pratiche utili in fase di progetto (convertitori, filtri, protezioni).
\end{abstract}

\tableofcontents

%======================================================================
\section{Dati di targa e ipotesi}

I dati considerati sono quelli minimi di targa:

\begin{center}
\begin{tabular}{llp{7.6cm}}
\toprule
Simbolo & Unità & Significato \\
\midrule
$S_n$      & VA & Potenza apparente nominale (trifase, somma delle tre fasi) \\
$U_n$      & V  & Tensione nominale \emph{concatenata} (lato considerato) \\
$I_n$      & A  & Corrente nominale \emph{di linea} (lato considerato) \\
$u_{cc}$   & \% & Tensione di cortocircuito percentuale \\
$P_{cc}$   & W  & Perdite di cortocircuito (nel rame) a corrente nominale --- opzionale \\
$f$        & Hz & Frequenza nominale \\
\bottomrule
\end{tabular}
\end{center}

\paragraph{Nota su $P_{nom}$.} Sulla targa dei trasformatori di potenza il dato
è la potenza \emph{apparente} $S_n$ in kVA. Se il dato disponibile è una potenza
\emph{attiva} $P_n$ con fattore di potenza $\cos\varphi$, si usa
\begin{equation}
S_n = \frac{P_n}{\cos\varphi}.
\end{equation}
Tutte le grandezze che seguono si basano su $S_n$: usare $P_n$ al suo posto porta a
sottostimare l'impedenza di base e quindi l'induttanza.

\paragraph{Ipotesi.} Si assume funzionamento simmetrico ed equilibrato, per cui il
trasformatore si rappresenta con il circuito monofase equivalente \emph{a stella},
indipendentemente dal collegamento fisico degli avvolgimenti. Le tre grandezze di targa
sono legate da
\begin{equation}
S_n = \sqrt{3}\, U_n I_n
\qquad\Longrightarrow\qquad
I_n = \frac{S_n}{\sqrt{3}\,U_n}.
\label{eq:Sn}
\end{equation}
$I_n$ è quindi ridondante: serve come verifica di coerenza dei dati.

%======================================================================
\section{Impedenza di base e sistema per unità}

Nel circuito monofase equivalente a stella, la tensione di fase è $U_n/\sqrt{3}$ e la
corrente di fase coincide con quella di linea $I_n$. L'impedenza di base è quindi
\begin{equation}
Z_b \;=\; \frac{U_n/\sqrt{3}}{I_n} \;=\; \frac{U_n}{\sqrt{3}\,I_n} \;=\; \frac{U_n^2}{S_n}.
\label{eq:Zb}
\end{equation}
Le tre espressioni sono equivalenti grazie alla \eqref{eq:Sn}. L'ultima è la più comoda:
richiede solo tensione concatenata e potenza apparente.

L'induttanza di base, utile per esprimere risultati in p.u., vale
\begin{equation}
L_b = \frac{Z_b}{\omega} = \frac{U_n^2}{2\pi f\, S_n}, \qquad \omega = 2\pi f .
\end{equation}

%======================================================================
\section{Significato fisico di $u_{cc}$}

$u_{cc}$ è definita dalla prova di cortocircuito: si cortocircuitano i morsetti di un
lato e si alimenta l'altro con una tensione ridotta, aumentata fino a far circolare la
corrente nominale. Il rapporto fra questa tensione e quella nominale è $u_{cc}$.

Poiché in tali condizioni la tensione applicata è piccola, il flusso nel nucleo è piccolo
e il ramo magnetizzante ($L_\mu$, $R_{fe}$) è trascurabile: il circuito si riduce alla
sola impedenza serie. Ne segue direttamente
\begin{equation}
\boxed{\;Z_{cc} \;=\; \frac{u_{cc}}{100}\; Z_b \;=\; \frac{u_{cc}}{100}\cdot\frac{U_n^2}{S_n}
\;=\;\frac{u_{cc}}{100}\cdot\frac{U_n}{\sqrt{3}\,I_n}\;}
\label{eq:Zcc}
\end{equation}
ovvero, in per unità, $z_{cc} = u_{cc}/100$.

$Z_{cc}$ è l'impedenza \emph{totale} di cortocircuito, cioè la serie di primario e
secondario riportati allo stesso lato. È la grandezza che conta per: corrente di guasto,
caduta di tensione sotto carico, $\mathrm{d}i/\mathrm{d}t$ imposto a un convertitore,
commutazione di un raddrizzatore.

%======================================================================
\section{Separazione fra resistenza e reattanza}

\subsection{Caso 1: $P_{cc}$ noto (preferibile)}

Le perdite di cortocircuito sono perdite Joule sulle tre fasi a corrente nominale:
\begin{equation}
P_{cc} = 3\,R_{cc}\,I_n^2
\qquad\Longrightarrow\qquad
R_{cc} = \frac{P_{cc}}{3 I_n^2},
\label{eq:Rcc}
\end{equation}
oppure, in forma p.u. (più robusta perché indipendente dal lato considerato),
\begin{equation}
u_R = \frac{P_{cc}}{S_n}\cdot 100 \;[\%],
\qquad R_{cc} = \frac{u_R}{100} Z_b .
\end{equation}
La reattanza si ottiene per differenza pitagorica:
\begin{equation}
\boxed{\;X_{cc} = \sqrt{Z_{cc}^2 - R_{cc}^2}
= \frac{Z_b}{100}\sqrt{u_{cc}^2 - u_R^2}\;}
\label{eq:Xcc}
\end{equation}

\subsection{Caso 2: $P_{cc}$ non noto}

Si può approssimare $X_{cc}\simeq Z_{cc}$. L'errore commesso è
\begin{equation}
\frac{X_{cc}}{Z_{cc}} = \sqrt{1-\left(\frac{R_{cc}}{X_{cc}}\right)^2\!\!\Big/\!\left(1+\left(\frac{R_{cc}}{X_{cc}}\right)^2\right)}
\;\simeq\; 1-\frac{1}{2}\left(\frac{R_{cc}}{Z_{cc}}\right)^2 ,
\end{equation}
quindi inferiore all'$1\%$ se $R_{cc}/Z_{cc}<0{,}14$. Valori tipici di $u_R$:

\begin{center}
\begin{tabular}{lcc}
\toprule
Taglia & $u_R$ tipico [\%] & $u_{cc}$ tipico [\%] \\
\midrule
$\le 100$ kVA (distribuzione)      & $1{,}5\div3$   & $4$ \\
$160\div1000$ kVA (distribuzione)  & $1{,}0\div1{,}5$ & $4\div6$ \\
$1\div10$ MVA                      & $0{,}5\div1{,}0$ & $6\div8$ \\
$>10$ MVA (potenza)                & $0{,}2\div0{,}5$ & $8\div16$ \\
\bottomrule
\end{tabular}
\end{center}

Sulle macchine grandi l'approssimazione $X_{cc}\simeq Z_{cc}$ è quindi eccellente; sui
piccoli trasformatori di distribuzione introduce qualche punto percentuale di errore.

%======================================================================
\section{Induttanza di dispersione totale}

Dalla reattanza si ricava immediatamente
\begin{equation}
\boxed{\;L_{cc} = \frac{X_{cc}}{2\pi f}
= \frac{1}{2\pi f}\cdot\frac{u_{cc}}{100}\cdot\frac{U_n^2}{S_n}\cdot
\sqrt{1-\left(\frac{u_R}{u_{cc}}\right)^{2}}\;}
\label{eq:Lcc}
\end{equation}
$L_{cc}$ è l'induttanza di dispersione \emph{totale per fase}, riferita al lato di cui si
sono usati $U_n$ e $S_n$, nel circuito equivalente a stella.

Formula compatta trascurando la resistenza:
\begin{equation}
L_{cc} \simeq \frac{u_{cc}}{100}\cdot\frac{U_n^2}{2\pi f\, S_n}
\qquad\text{(errore }<1\%\text{ se } u_R/u_{cc}<0{,}14\text{)}.
\end{equation}

%======================================================================
\section{Ripartizione fra primario e secondario}

I dati di targa forniscono solo la \emph{somma} $L_{cc}=L_{\sigma 1}+L'_{\sigma 2}$: la
prova di cortocircuito non permette in alcun modo di distinguere i due contributi. La
separazione richiede o dati costruttivi, o misure aggiuntive, o una FEM.

La convenzione usuale, valida con buona approssimazione per avvolgimenti concentrici di
uguale altezza, è la ripartizione simmetrica \emph{in p.u.}:
\begin{equation}
L_{\sigma 1} = L'_{\sigma 2} = \frac{L_{cc}}{2}.
\end{equation}

Riportando al secondario con il rapporto di trasformazione del circuito equivalente
(rapporto fra tensioni \emph{concatenate})
\begin{equation}
a = \frac{U_{n1}}{U_{n2}},
\qquad
L_{\sigma 2} = \frac{L'_{\sigma 2}}{a^{2}} = \frac{L_{cc}}{2\,a^{2}} .
\end{equation}

In alternativa, e in modo del tutto equivalente, si può ripartire l'errore ricalcolando
tutto direttamente sul lato secondario con la \eqref{eq:Lcc} usando $U_{n2}$: il valore
p.u. $u_{cc}$ è invariante rispetto al lato.

\paragraph{Attenzione.} Il rapporto $a$ del circuito equivalente \emph{non} coincide in
generale con il rapporto spire reale $N_1/N_2$. Per un Dyn11, ad esempio,
$a = U_{n1}/U_{n2}$ mentre $N_1/N_2 = U_{n1}/(U_{n2}/\sqrt{3}) = \sqrt{3}\,a$.

%======================================================================
\section{Collegamento a triangolo: fattore 3}
\label{sec:delta}

Tutte le formule precedenti danno l'induttanza del circuito equivalente \emph{a stella}.
Se l'avvolgimento è fisicamente a triangolo e serve l'induttanza dell'avvolgimento reale
(ad esempio per un modello circuitale dettagliato, per una simulazione dei transitori di
commutazione o per il dimensionamento di uno snubber), vale la trasformazione
stella--triangolo:
\begin{equation}
\boxed{\;L_{\sigma,\Delta} = 3\, L_{\sigma,Y\text{-eq}}\;}
\end{equation}

Analogamente, se si misura l'induttanza di cortocircuito con un LCR fra due morsetti di
linea (l'altro lato cortocircuitato), su un avvolgimento a stella si legge
\begin{equation}
L_{cc,\,LL} = 2\,L_{cc,\,fase}\, ,
\end{equation}
mentre su un triangolo si legge il parallelo fra un ramo e la serie degli altri due:
$L_{cc,LL} = \tfrac{2}{3}L_{\sigma,\Delta} = 2 L_{\sigma,Y\text{-eq}}$, cioè lo stesso
valore --- come deve essere, perché il comportamento ai morsetti non dipende dal
collegamento interno.

%======================================================================
\section{Esempio numerico}

Trasformatore di distribuzione in olio:
\begin{center}
\begin{tabular}{lll}
\toprule
$S_n = 630$ kVA & $U_{n1}=20$ kV & $U_{n2}=400$ V \\
$u_{cc}=6\%$ & $P_{cc}=6500$ W & $f=50$ Hz \\
Collegamento Dyn11 & & \\
\bottomrule
\end{tabular}
\end{center}

\subsection*{Correnti nominali}
\begin{align}
I_{n1} &= \frac{630\cdot 10^3}{\sqrt{3}\cdot 20\cdot 10^3} = 18{,}19\un{A}\\
I_{n2} &= \frac{630\cdot 10^3}{\sqrt{3}\cdot 400} = 909{,}3\un{A}
\end{align}

\subsection*{Lato MT (primario)}
\begin{align}
Z_{b1} &= \frac{U_{n1}^2}{S_n} = \frac{(20\cdot10^3)^2}{630\cdot10^3} = 634{,}9\un{\Omega}\\
Z_{cc1} &= 0{,}06\cdot 634{,}9 = 38{,}10\un{\Omega}\\
u_R &= \frac{6500}{630\cdot10^3}\cdot100 = 1{,}032\,\% \;\Rightarrow\;
R_{cc1} = 0{,}01032\cdot634{,}9 = 6{,}55\un{\Omega}\\
X_{cc1} &= \sqrt{38{,}10^2-6{,}55^2} = 37{,}53\un{\Omega}
\qquad (u_X = 5{,}91\,\%)\\
L_{cc1} &= \frac{37{,}53}{2\pi\cdot 50} = 119{,}5\un{mH}
\end{align}
Si noti che trascurare $R_{cc}$ avrebbe dato $L_{cc1}=121{,}3$ mH, con un errore del
$+1{,}5\%$: accettabile nella maggior parte dei casi.

\subsection*{Ripartizione e riporto}
\begin{align}
L_{\sigma 1} = L'_{\sigma 2} &= \frac{119{,}5}{2} = 59{,}7\un{mH}\\
a &= \frac{20000}{400} = 50\\
L_{\sigma 2} &= \frac{59{,}7\cdot10^{-3}}{50^2} = 23{,}9\un{\mu H}
\end{align}
Verifica dal lato BT: $Z_{b2}=400^2/630\cdot10^3 = 0{,}2540\ \Omega$,
$X_{cc2}=0{,}0591\cdot0{,}2540 = 15{,}01\ \mathrm{m\Omega}$,
$L_{cc2}=47{,}8\ \mu\mathrm{H}$, metà $=23{,}9\ \mu$H. \checkmark

\subsection*{Avvolgimento MT fisico (triangolo)}
Il primario è a triangolo, quindi l'induttanza di dispersione dell'avvolgimento reale di
una colonna, riferita al primario, è
\begin{equation}
L_{cc,\Delta} = 3\cdot 119{,}5 = 358{,}6\un{mH},
\qquad L_{\sigma 1,\Delta} = 3\cdot 59{,}7 = 179{,}3\un{mH}.
\end{equation}

\subsection*{Corrente di cortocircuito}
Trascurando l'impedenza di rete a monte (rete di potenza infinita):
\begin{equation}
I_{cc} = \frac{I_n}{u_{cc}/100} = \frac{I_n}{0{,}06} = 16{,}7\,I_n
\quad\Rightarrow\quad
I_{cc1}=303\un{A},\qquad I_{cc2}=15{,}2\un{kA}.
\end{equation}
Il valore di picco asimmetrico può arrivare a $\sqrt2\,\kappa\,I_{cc}$ con
$\kappa\simeq1{,}8$ per $X/R\simeq 5{,}7$, cioè circa $38$ kA di picco lato BT.

%======================================================================
\section{Verifica geometrica indipendente: formula di Rogowski}

L'induttanza di dispersione è un parametro essenzialmente \emph{geometrico}: il flusso
disperso si richiude nel canale fra gli avvolgimenti, cioè in aria/olio. Per avvolgimenti
concentrici della stessa altezza vale la classica formula
\begin{equation}
\boxed{\;L_{\sigma} = \mu_0 N_1^2 \,\frac{\pi D_m}{h_w}\,
\left(a_{12} + \frac{a_1+a_2}{3}\right) K_R \;}
\label{eq:rogowski}
\end{equation}
dove $N_1$ è il numero di spire dell'avvolgimento a cui si riferisce il risultato,
$D_m$ il diametro medio del canale, $h_w$ l'altezza degli avvolgimenti, $a_{12}$ lo
spessore radiale del canale fra gli avvolgimenti, $a_1,a_2$ gli spessori radiali dei due
avvolgimenti, e
\begin{equation}
K_R = 1 - \frac{a_{12}+a_1+a_2}{\pi h_w}
\end{equation}
è il fattore di Rogowski, che corregge l'ipotesi di campo puramente assiale ai bordi.

\paragraph{Applicazione all'esempio.} Con dimensioni plausibili per un 630 kVA
($N_1 = 1900$ spire lato MT a triangolo, $D_m = 0{,}28$ m, $h_w = 0{,}45$ m,
$a_{12}=20$ mm, $a_1=25$ mm, $a_2=30$ mm):
\begin{align}
a_{12}+\frac{a_1+a_2}{3} &= 0{,}020 + 0{,}0183 = 0{,}0383\un{m}\\
K_R &= 1-\frac{0{,}075}{\pi\cdot0{,}45} = 0{,}947\\
L_\sigma &= 4\pi\!\cdot\!10^{-7}\cdot 1900^2\cdot\frac{\pi\cdot0{,}28}{0{,}45}\cdot
0{,}0383\cdot 0{,}947 = 0{,}322\un{H}
\end{align}
da confrontare con i $358{,}6$ mH ricavati dalla targa (avvolgimento a triangolo,
Sez.~\ref{sec:delta}): scarto di circa il $10\%$, del tutto ragionevole avendo
\emph{ipotizzato} la geometria. Il confronto ha valore soprattutto come controllo di
ordine di grandezza e come verifica che si stia confrontando la stessa grandezza (stella
equivalente contro avvolgimento fisico).

Dalla \eqref{eq:rogowski} discendono alcune conseguenze progettuali utili:
\begin{itemize}
  \item $L_\sigma \propto N^2$ e $\propto 1/h_w$: avvolgimenti alti e sottili
        riducono la dispersione;
  \item $L_\sigma \propto a_{12}$: il canale di isolamento fra MT e BT domina, ed è
        imposto dalla tenuta dielettrica. Da qui il legame fra classe di tensione e
        $u_{cc}$ tipica (tabella precedente);
  \item per aumentare deliberatamente $L_\sigma$ (trasformatori per DAB, LLC, forni ad
        arco, saldatrici) si allarga il canale o si inseriscono shunt magnetici;
  \item per ridurla si interlacciano gli avvolgimenti (interleaving), che riduce
        $L_\sigma$ circa con $1/p^2$ dove $p$ è il numero di sezioni.
\end{itemize}

%======================================================================
\section{Note pratiche e casi particolari}

\paragraph{L'induttanza di dispersione non satura.} Il flusso disperso si richiude
prevalentemente in aria/olio, quindi $L_{cc}$ è sostanzialmente costante anche in
cortocircuito franco, quando il nucleo è invece ampiamente saturo. Questa è la ragione per
cui il modello a impedenza costante resta valido per il calcolo delle correnti di guasto.

\paragraph{Dipendenza dalla frequenza.} A 50/60 Hz $L_{cc}$ è praticamente costante. Sopra
qualche kHz (armoniche di convertitore, transitori di manovra) l'effetto pelle e di
prossimità riducono leggermente $L_{cc}$ (il flusso viene espulso dai conduttori) e
aumentano fortemente $R_{cc}$. Per studi EMI/transitori occorre inoltre includere le
capacità parassite: la frequenza di risonanza propria dell'avvolgimento è spesso nel range
$10\div500$ kHz.

\paragraph{Variatore di prese.} $u_{cc}$ è riferita alla presa principale. Sulle prese
estreme può variare di $\pm10\div20\%$ (ampiezza che dipende dalla posizione fisica delle
spire di regolazione). Se il dato è critico, va richiesto il valore di $u_{cc}$ su tutte
le prese.

\paragraph{Temperatura.} Secondo IEC 60076-1 le perdite si riferiscono a una temperatura
convenzionale (75\,°C o 115\,°C secondo la classe). Solo $R_{cc}$ dipende dalla
temperatura:
\begin{equation}
R_{cc}(\vartheta) = R_{cc}(\vartheta_0)\,\frac{235+\vartheta}{235+\vartheta_0}
\quad\text{(rame)},
\end{equation}
mentre $X_{cc}$ e quindi $L_{cc}$ ne sono indipendenti. È un ulteriore motivo per
preferire il calcolo di $L$ via $X_{cc}$ anziché via $Z_{cc}$.

\paragraph{Sequenza omopolare.} $L_{cc}$ calcolata qui è quella di \emph{sequenza
diretta}. L'impedenza omopolare $Z_0$ dipende dal collegamento e dal tipo di nucleo:
\begin{itemize}
  \item Dyn: $Z_0 \simeq (0{,}85\div1{,}0)\,Z_{cc}$ (il triangolo offre un percorso alle
        correnti di circolazione);
  \item YNyn su nucleo a 3 colonne: $Z_0 \simeq (0{,}3\div1{,}0)\,Z_{cc}$, fortemente
        dipendente dal ritorno di flusso in cassa;
  \item YNyn su banco di tre monofase o nucleo a 5 colonne: $Z_0$ molto elevata
        ($10\div100$ p.u.).
\end{itemize}
Per i calcoli di guasto a terra il dato va richiesto esplicitamente al costruttore.

\paragraph{Trasformatori a tre avvolgimenti.} Con tre avvolgimenti la targa riporta tre
tensioni di cortocircuito $u_{12}, u_{13}, u_{23}$ (ciascuna riferita alla propria base di
potenza, da riportare a una base comune). Il modello è una stella di tre impedenze:
\begin{equation}
Z_1 = \tfrac{1}{2}(Z_{12}+Z_{13}-Z_{23}),\quad
Z_2 = \tfrac{1}{2}(Z_{12}+Z_{23}-Z_{13}),\quad
Z_3 = \tfrac{1}{2}(Z_{13}+Z_{23}-Z_{12}).
\end{equation}
Uno dei tre può risultare negativo: è un artificio matematico del modello, non un errore.

\paragraph{Uso in elettronica di potenza.} Nei convertitori grid-tied la dispersione del
trasformatore di accoppiamento è spesso usata come induttanza di filtro lato rete. Con
$u_{cc}=6\%$ si ha per definizione $x = 0{,}06$ p.u.: è già un valore adeguato per un
filtro L (tipicamente $0{,}05\div0{,}15$ p.u.), ma va verificato il ripple:
\begin{equation}
\Delta i_{pp} \simeq \frac{U_{dc}}{8 f_{sw} L_{cc}} \quad\text{(PWM a due livelli, per fase)} .
\end{equation}
Nei DAB/LLC, invece, $L_{cc}$ è l'elemento di trasferimento di potenza e va progettata,
non subita: in quel caso il dato di targa serve solo come punto di partenza.

\paragraph{Misura diretta.} In laboratorio, senza dati di targa:
\begin{enumerate}
  \item cortocircuitare saldamente il lato BT (barre di sezione adeguata: la resistenza
        del cortocircuito entra direttamente nella misura);
  \item alimentare il lato MT a tensione ridotta, sinusoidale, alla frequenza nominale;
  \item misurare $U$, $I$, $P$ trifase e ricavare
        $Z_{cc}=U/(\sqrt3 I)$, $R_{cc}=P/(3I^2)$, $X_{cc}=\sqrt{Z_{cc}^2-R_{cc}^2}$;
  \item se possibile portarsi vicino a $I_n$: a corrente molto ridotta il fattore di
        potenza è basso ($\cos\varphi \simeq 0{,}17$) e la misura di $P$ è critica.
\end{enumerate}
Un LCR meter a 50--100 Hz è accettabile per una stima rapida, ma sottostima
sistematicamente $R_{cc}$ (misura a corrente bassissima, senza effetto pelle
significativo).

%======================================================================
\section{Riepilogo delle formule}

\begin{center}
\begin{tabular}{lll}
\toprule
Grandezza & Formula & Esempio (630 kVA, lato MT)\\
\midrule
Corrente nominale     & $I_n = S_n/(\sqrt3 U_n)$              & $18{,}19$ A \\
Impedenza di base     & $Z_b = U_n^2/S_n$                      & $634{,}9\ \Omega$ \\
Impedenza di cc       & $Z_{cc} = (u_{cc}/100) Z_b$            & $38{,}10\ \Omega$ \\
Resistenza di cc      & $R_{cc} = P_{cc}/(3I_n^2)$             & $6{,}55\ \Omega$ \\
Reattanza di cc       & $X_{cc} = \sqrt{Z_{cc}^2-R_{cc}^2}$    & $37{,}53\ \Omega$ \\
Induttanza totale     & $L_{cc} = X_{cc}/(2\pi f)$             & $119{,}5$ mH \\
Per avvolgimento      & $L_{\sigma 1}=L'_{\sigma 2}=L_{cc}/2$  & $59{,}7$ mH \\
Riporto al secondario & $L_{\sigma 2} = L'_{\sigma 2}/a^2$     & $23{,}9\ \mu$H \\
Avvolgimento a $\Delta$ & $L_\Delta = 3 L_{Y\text{-eq}}$       & $179{,}3$ mH \\
\bottomrule
\end{tabular}
\end{center}

Formula unica, trascurando la resistenza:
\begin{equation}
\boxed{\;L_{cc} \simeq \frac{u_{cc}}{100}\cdot \frac{U_n^2}{2\pi f\,S_n}\;}
\end{equation}

\end{document}
